\documentclass[11pt]{ctexart}
\usepackage[a4paper,margin=1in]{geometry}
\usepackage{graphicx}
\usepackage{xcolor}
\usepackage{hyperref}
\definecolor{refgray}{gray}{0.45}
\title{\textbf{市场最新观点汇总}\\[0.4em]{\large AI与机器人主导的结构性分化成为全球市场主线，权益风险溢价压缩与加息风险回归构成主要宏观约束。}}
\author{KC桌面 自动生成}
\date{260521}
\begin{document}
\maketitle
\section*{一页摘要}
\begin{itemize}
  \item 全球投资者风险偏好框架发生结构性重置，对AI踏空的恐惧已取代地缘政治成为首要‘存在性威胁’，AI资本支出被视为驱动市场表现的压倒性力量。
  \item 中国在机器人领域已建立系统性优势，2025年全球85\%的人形机器人部署发生在中国，供应链的‘基础设施化’是市场真正低估的变量。
  \item 美国权益风险溢价已跌破2007年低位，债券相对于股票的吸引力正在回升，市场抵御利率进一步上行的‘安全垫’极为有限。
  \item 美联储加息风险正在重现，两年期美债收益率与联邦基金利率的利差信号指向未来一年内加息的可能性，这是当前最被低估的宏观风险。
  \item 亚洲出口复苏并非普涨，而是由AI驱动的结构性分化，半导体和AI相关链条是主要驱动力，台湾和韩国显著跑赢区域平均水平。
\end{itemize}
\section{宏观与利率}
\textbf{市场对降息叙事的依赖正在被打破，加息风险与权益风险溢价的历史性压缩构成当前资产定价的核心约束。}

\begin{itemize}
  \item 美国消费者信心调查中的净利率预期指标已明确转向‘未来一年利率将更高’，同时两年期美债收益率与联邦基金利率的利差升至约45个基点，这一利差在历史上领先联邦基金利率变化约十一个月，相关性高达66\%。
  \item 标普500的权益风险溢价已降至2.2\%，不仅低于3.1\%的长期历史均值约90个基点，更突破了2007年的前低，创下雷曼危机以来的最低纪录，这意味着权益市场对利率上行的缓冲空间已极其有限。
  \item 全球综合债券指数收益率已逼近4\%，美国10年期实际收益率同步走高，但真正值得关注的不是利率绝对水平，而是权益与债券之间的相对定价关系已出现历史级别的收窄。
  \item A股W\&W情绪指标从‘非常看多’的极端区间回落至‘温和看多’，月度均值29仍处于看多区间，但市场向上弹性的确定性在下降，结构选择比方向判断更重要。
  \item 人民币中间价定价模型的误差长期偏移暗示政策框架的隐性调整，逆周期因子等政策工具的‘使用强度’正在发生变化，汇率形成机制可能正从被动维稳转向主动校准。
\end{itemize}
\begin{figure}[htbp]
\centering
\includegraphics[width=0.88\linewidth]{figures/fig_047.jpg}
\caption{Figure 1 - Figure 1: Rising spread between 2y Treasury yield and fed funds rate pointing to rate hike risk}
\end{figure}
\begin{figure}[htbp]
\centering
\includegraphics[width=0.88\linewidth]{figures/fig_072.jpg}
\caption{Figure 1 - Figure 1: Our model based S\&P500 Equity Discount Rate vs. the real 10y UST yield}
\end{figure}
\begin{figure}[htbp]
\centering
\includegraphics[width=0.88\linewidth]{figures/fig_073.jpg}
\caption{Figure 1 - Figure 2: Equity Risk Premium proxied by the difference between the S\&P500 Equity Discount Rate and the real 10y UST yield}
\end{figure}
\begin{figure}[htbp]
\centering
\includegraphics[width=0.88\linewidth]{figures/fig_031.jpg}
\caption{Exhibit 1 - Exhibit 1: BofA China A-share W\&W Indicator in May 2026 The latest daily reading at 37 is "marginally bullish". The monthly average at 29 is bullish BofA China A-share Wax \& Wane Indicator}
\end{figure}
\begin{figure}[htbp]
\centering
\includegraphics[width=0.88\linewidth]{figures/fig_032.jpg}
\caption{Exhibit 2 - Exhibit 2: W\&W Indicator vs percentage change in CSI300 in the following 1mth (Back-tested) since Jan-2015 W\&W Indicator has been fairly good in predicting the CSI300 1mth performance}
\end{figure}
{\small\color{refgray}\textbf{References:} [R005] 市场真正低估的不是通胀，而是美联储加息风险的重现; [R009] 市场真正低估的不是地缘风险，而是权益风险溢价已跌破2007年低位; [R003] A股信号已从“非常看多”转向“温和看多”，真正的考验在于结构而非方向; [R012] 人民币中间价定价模型正在发出一个被忽视的信号}

\section{AI/算力与机器人}
\textbf{AI正在从训练侧向推理和边缘侧加速渗透，Agentic AI的兴起正在重构CPU服务器需求，而机器人供应链的‘基础设施化’是中国主导全球机器人产业的核心变量。}

\begin{itemize}
  \item 中国在机器人领域的主导地位并非单点技术突破，而是由政策、供应链、人口结构和创新体系共同编织的系统性优势，2025年全球人形机器人部署量的85\%发生在中国，中国制造商宇树科技和智元机器人合计贡献超过70\%的装机量。
  \item 人形机器人正在从‘实验室验证’跨过‘小批量交付’的早期商业化门槛，2025年全球安装量从约2000台跃升至约15000台，供应链的规模效应开始启动，核心零部件成本曲线将出现非线性下降。
  \item Agentic AI的推理过程涉及多步骤任务执行、工具调用和大量数据移动，对通用计算和编排能力的需求远高于对纯GPU算力的依赖，CPU服务器市场总规模到2030年可能扩张至当前的2.5到4.5倍。
  \item AI硬件拥挤度虽处历史高位，但过去三个季度并未进一步恶化，资金正在从‘买赛道’转向‘挑环节’，电子板块配置比例从峰值回落而电信板块创下新高，结构分化正在创造alpha机会。
  \item AI供应链的瓶颈正在从GPU本身向上游的被动元器件、封装材料和功率芯片迁移，Murata将AI服务器MLCC需求CAGR从30\%大幅上调至80\%，日本电子元器件企业正在获得结构性增长的定价权。
  \item Rubin机架单机架成本从Blackwell的约400万美元跃升至约780万美元，价值在供应链中重新分配，PCB、MLCC等被动元器件的单位价值量跃升正在重塑竞争壁垒。
\end{itemize}
\begin{figure}[htbp]
\centering
\includegraphics[width=0.88\linewidth]{figures/fig_004.jpg}
\caption{FIGURE 4 - Note: 2026 data shows the consensus estimate. FIGURE 4. Chinese companies delivered 85\% of global humanoid installations}
\end{figure}
\begin{figure}[htbp]
\centering
\includegraphics[width=0.88\linewidth]{figures/fig_005.jpg}
\caption{Figure 5 - \# Robotics innovation has shifted to China \# Our analysis reveals that China accounts for 70\% of global robotics inventions since 2000 (Figure 5). For example, in 2023 – the latest year with largely complete data – China recorded over 30k robotics patent inven}
\end{figure}
\begin{figure}[htbp]
\centering
\includegraphics[width=0.88\linewidth]{figures/fig_012.jpg}
\caption{Figure 15 - With nearly 13k humanoid robots deployed in China by end-2025, early scaling is already generating powerful second-order effects. Larger deployments are producing growing volumes of real-world operating data, which directly improve humanoid robot models – help}
\end{figure}
\begin{figure}[htbp]
\centering
\includegraphics[width=0.88\linewidth]{figures/fig_014.jpg}
\caption{Figure 17 - \# Humanoids are (electric) cars in miniature At this early stage, most bottom-up forecasts for humanoid robots rest on highly uncertain assumptions around unit costs, utilisation, returns on investment and regulatory acceptance. We assess humanoid manufacturin}
\end{figure}
\begin{figure}[htbp]
\centering
\includegraphics[width=0.88\linewidth]{figures/fig_374.jpg}
\caption{EXHIBIT 1 - EXHIBIT 1: Chatbot inference workflow \# Chatbot Inference Single-pass: prompt in \$\textbackslash{}rightarrow\$ answer out ```mermaid graph LR}
\end{figure}
\begin{figure}[htbp]
\centering
\includegraphics[width=0.88\linewidth]{figures/fig_375.jpg}
\caption{EXHIBIT 2 - EXHIBIT 2: Agentic inference workflow ```mermaid graph TD}
\end{figure}
\begin{figure}[htbp]
\centering
\includegraphics[width=0.88\linewidth]{figures/fig_376.jpg}
\caption{EXHIBIT 3 - EXHIBIT 3: CPU Vendor Comments on Agentic AI driving CPU demand growth EXHIBIT 4: 2.5x-4.5x Traditional server TAM expansion by 2030 CPU server TAM BY 2030}
\end{figure}
\begin{figure}[htbp]
\centering
\includegraphics[width=0.88\linewidth]{figures/fig_086.jpg}
\caption{Figure 2 - Figure 1: Onshore mutual fund holdings in A share electronics}
\end{figure}
\begin{figure}[htbp]
\centering
\includegraphics[width=0.88\linewidth]{figures/fig_087.jpg}
\caption{Figure 2 - Figure 2: Onshore mutual fund holdings in A share telecom and electronics}
\end{figure}
\begin{figure}[htbp]
\centering
\includegraphics[width=0.88\linewidth]{figures/fig_088.jpg}
\caption{Figure 3 - Figure 3: Crowding scores vs performance of popular China AI names (ex-internet)}
\end{figure}
\begin{figure}[htbp]
\centering
\includegraphics[width=0.88\linewidth]{figures/fig_351.jpg}
\caption{Exhibit 1 - Exhibit 1: Owing to recent increase in memory prices, memory will become 25\%+ of the rack BOM for Rubin Exhibit 2: We estimate the new upcoming VR200 rack ASP to be \textbackslash{}\textasciitilde{}US\textbackslash{}\$7.8mn Nvidia NVL72 Rack ASP (US\textbackslash{}\$)}
\end{figure}
\begin{figure}[htbp]
\centering
\includegraphics[width=0.88\linewidth]{figures/fig_352.jpg}
\caption{Exhibit 1 - Exhibit 1: Owing to recent increase in memory prices, memory will become 25\%+ of the rack BOM for Rubin Exhibit 2: We estimate the new upcoming VR200 rack ASP to be \textbackslash{}\textasciitilde{}US\textbackslash{}\$7.8mn Nvidia NVL72 Rack ASP (US\textbackslash{}\$)}
\end{figure}
{\small\color{refgray}\textbf{References:} [R001] 机器人十年属于中国，但市场真正低估的不是产能，而是供应链的“基础设施化”; [R002] 人形机器人真正改变的不是就业，而是宏观增长方程式的分母; [R042] 市场真正低估的不是AI训练，而是Agentic AI对传统服务器的需求重构; [R011] AI硬件拥挤度已至历史高位，但市场真正低估的不是拥挤本身，而是拥挤背后的结构分化; [R045] AI供应链的真正拐点：不是GPU本身，而是围绕它的电子元器件正在经历一次供给侧的定价重构; [R037] 市场低估了Rubin机架带来的价值重分配，而非仅仅是算力升级; [R044] 市场低估了AI基础设施的“Token工厂”模式，而非算力租赁本身; [R040] 半导体分销的AI红利，市场仍低估了规模效应的复利}

\section{能源与大宗商品}
\textbf{供给端的刚性约束正在为部分大宗商品建立新的价格底部，油价的‘粘性’和中间环节的供给缺口是市场真正低估的结构性变量。}

\begin{itemize}
  \item 地缘政治风险指数（GPR）与油价的实证分析表明，即使地缘风险迅速消散，油价的下跌路径远比历史经验所暗示的要慢，布伦特原油的模型‘公允价值’约在每桶120美元，远高于当前期货曲线。
  \item 日本原油进口价格在两个月内飙升57\%，而进口量却暴跌63.7\%，这一‘价格飞涨、数量骤降’的组合揭示了全球供应链重构背景下能源安全的深层变化。
  \item 2026年磷酸铁（FePO4）将出现供给短缺，新增产能有限且爬坡周期长，预计价格还有2000-3000元/吨的上涨空间，拥有垂直一体化能力或采用不同工艺路线的正极厂商将受益。
  \item 中国4月经济数据全面走弱，但供给端的约束正在为部分大宗商品建立新的价格底部，铝的结构性短缺逻辑仍在，钢铁利润端因成本推动出现回暖。
  \item 巴西糖市当前供给端的放松是以牺牲未来供给为代价的，生产成本高于16美分/磅而近月合约仅15美分/磅，蔗农利润被压缩至亏损线以下，2027/28年度全球平衡表可能因此收紧。
\end{itemize}
\begin{figure}[htbp]
\centering
\includegraphics[width=0.88\linewidth]{figures/fig_121.jpg}
\caption{EXHIBIT 1 - EXHIBIT 1: Across over 90 quarters of data, the model explains \textbackslash{}\textasciitilde{}70\% of the variation in real oil prices. \# THE FOUR DRIVERS THAT ACTUALLY MATTER Ranking by economic impact — how much a one-standard-deviation move in each variabl}
\end{figure}
\begin{figure}[htbp]
\centering
\includegraphics[width=0.88\linewidth]{figures/fig_122.jpg}
\caption{EXHIBIT 2 - EXHIBIT 2: GPR can spike for non-oil reasons (9-11 for example) and oil price can spike without a GPR signal (the 2008 spike) but generally oil conflicts drive up both GPR and price, driving supply shocks}
\end{figure}
\begin{figure}[htbp]
\centering
\includegraphics[width=0.88\linewidth]{figures/fig_125.jpg}
\caption{EXHIBIT 5 - EXHIBIT 5: Brent prices exhibit greater persistence than the GPR index during past supply shock episodes Post-Crisis Persistence (Supply Shocks Only)}
\end{figure}
\begin{figure}[htbp]
\centering
\includegraphics[width=0.88\linewidth]{figures/fig_128.jpg}
\caption{EXHIBIT 8 - EXHIBIT 8: While the GPR index has gradually normalized to the 200 level, Brent prices have shown greater stickiness}
\end{figure}
\begin{figure}[htbp]
\centering
\includegraphics[width=0.88\linewidth]{figures/fig_050.jpg}
\caption{Exhibit 3 - Exhibit 3: Japan's Imported Crude Oil and LNG Volume}
\end{figure}
\begin{figure}[htbp]
\centering
\includegraphics[width=0.88\linewidth]{figures/fig_051.jpg}
\caption{Exhibit 4 - Exhibit 4: Japan's Imported Oil Price and Dubai Oil Price}
\end{figure}
\begin{figure}[htbp]
\centering
\includegraphics[width=0.88\linewidth]{figures/fig_357.jpg}
\caption{Exhibit 7 - Exhibit 7: Rapidly rising ethanol inventories are pressuring ethanol prices, reducing immediate off-take capacity. Exhibit 8: Falling Ethanol and Normalizing Oil Reinforce Brazil Max-Sugar Expectations}
\end{figure}
\begin{figure}[htbp]
\centering
\includegraphics[width=0.88\linewidth]{figures/fig_358.jpg}
\caption{Exhibit 7 - Exhibit 7: Rapidly rising ethanol inventories are pressuring ethanol prices, reducing immediate off-take capacity. Exhibit 8: Falling Ethanol and Normalizing Oil Reinforce Brazil Max-Sugar Expectations}
\end{figure}
\begin{figure}[htbp]
\centering
\includegraphics[width=0.88\linewidth]{figures/fig_359.jpg}
\caption{Exhibit 1 - Exhibit 1: Sugar Stocks to use (\%) and NY Sugar Price (Usc/lb): We expect Global Stock-to-use ratio declining toward historical levels, supporting prices}
\end{figure}
{\small\color{refgray}\textbf{References:} [R013] 市场真正低估的不是地缘风险，而是油价的“粘性”; [R006] 日本贸易数据揭示的真正风险：不是出口放缓，而是能源供给的“价格-数量”错位; [R024] 电池材料行业真正被低估的不是需求增速，而是中间环节的供给缺口正在重新定义定价权; [R029] 市场低估的不是需求疲弱，而是供给端正在形成的定价权; [R038] 糖市真正的拐点不在巴西压榨量，而在被低估的供给反噬风险}

\section{地产与消费}
\textbf{中国房地产市场正在经历供给侧的主动收缩与需求端的K型分化，一线城市二手房价格企稳但信用传导时滞较长，消费领域物流成本通胀正在侵蚀利润表。}

\begin{itemize}
  \item 一线城市二手房价格连续第二个月全面正增长，四个城市方向统一，但市场真正在定价的是供给侧的持续收缩，新开工面积同比降幅从-17\%扩大至-27\%，房价企稳与供给收缩正在形成‘被动的再平衡’。
  \item 4月居民中长期贷款净减少3410亿元创历史新低，提前还贷规模仍在扩大，从市场企稳到信用扩张之间有一条漫长的传导链，居民加杠杆意愿仍然极低。
  \item 中国奢侈品市场呈现K型复苏，高端消费依旧强劲而中产阶层购买力仍然脆弱，品牌两极分化加剧，爱马仕、LV等强势品牌持续扩张而二、三线品牌面临客流和销售的双重压力。
  \item 物流成本通胀正在从3-4月开始实质性侵蚀企业利润，柴油价格较2025年均价已上涨11\%，啤酒、乳制品和部分饮料企业承受的冲击最为显著，净利润下行空间在低个位数到高个位数之间。
  \item 四月车市渠道库存减少了2.3万辆，供给端正在经历主动或被动的出清，合资与豪华品牌集体失速并非需求问题而是供给结构性问题，消费者正在用脚投票加速电动化转型。
\end{itemize}
\begin{figure}[htbp]
\centering
\includegraphics[width=0.88\linewidth]{figures/fig_316.jpg}
\caption{Figure 1 - Figure 1: Tier-1 cities - secondary home price index (NBS \& Centraline) with major nationwide easing/narrative change}
\end{figure}
\begin{figure}[htbp]
\centering
\includegraphics[width=0.88\linewidth]{figures/fig_317.jpg}
\caption{Figure 2 - Figure 2: Tier-1 cities - secondary home price M/M growth (NBS vs. Centraline)}
\end{figure}
\begin{figure}[htbp]
\centering
\includegraphics[width=0.88\linewidth]{figures/fig_322.jpg}
\caption{Figure 7 - Figure 7: NBS 70-city home price index - M/M growth}
\end{figure}
\begin{figure}[htbp]
\centering
\includegraphics[width=0.88\linewidth]{figures/fig_292.jpg}
\caption{Figure 1 - Figure 1: Household mid-long term new lending}
\end{figure}
\begin{figure}[htbp]
\centering
\includegraphics[width=0.88\linewidth]{figures/fig_293.jpg}
\caption{Figure 2 - Figure 2: Commercial mortgage rate and housing provident fund rate}
\end{figure}
\begin{figure}[htbp]
\centering
\includegraphics[width=0.88\linewidth]{figures/fig_294.jpg}
\caption{Figure 4 - Figure 4: MTD primary home sales by GFA in 30 city: tier-1 cities take lead with \$7\textbackslash{}\%\$ y-o-y Figure 5: MTD average secondary units sold in 10 cities +22\% y-o-y}
\end{figure}
\begin{figure}[htbp]
\centering
\includegraphics[width=0.88\linewidth]{figures/fig_146.jpg}
\caption{EXHIBIT 1 - EXHIBIT 1: Resilient demand at the higher end continues to support RevPAR outperformance within the luxury hotel segment. The middle tiers are likely most squeezed - on both a 2yr and 3yr basis, this group has even underperformed t}
\end{figure}
\begin{figure}[htbp]
\centering
\includegraphics[width=0.88\linewidth]{figures/fig_147.jpg}
\caption{EXHIBIT 2 - EXHIBIT 2: The recovery in property prices, which began in Sep-25, has accelerated since Jan-26, with week-on-week price growth averaging approximately +0.4\%}
\end{figure}
\begin{figure}[htbp]
\centering
\includegraphics[width=0.88\linewidth]{figures/fig_148.jpg}
\caption{EXHIBIT 3 - EXHIBIT 3: Improved consumer sentiment/ feel-good, driven by rising property and equity markets in Hong Kong, has translated into a buoyant retail environment, with particularly strong momentum in the watches and jewellery category}
\end{figure}
\begin{figure}[htbp]
\centering
\includegraphics[width=0.88\linewidth]{figures/fig_277.jpg}
\caption{Exhibit 1 - Exhibit 1: Industry's monthly inventory changes Inventory at channels decreased 23k units in April, after adjusting for export volumes}
\end{figure}
\begin{figure}[htbp]
\centering
\includegraphics[width=0.88\linewidth]{figures/fig_278.jpg}
\caption{Exhibit 2 - Exhibit 2: Monthly vehicle insurance data (vehicles manufactured in China only) April PV insurance registrations came in at 1.36mn units, down \$19\textbackslash{}\%\$ YoY/7\% MoM}
\end{figure}
{\small\color{refgray}\textbf{References:} [R030] 一线城市二手房价格企稳，但市场真正在定价的，是供给侧的持续收缩; [R027] 市场真正低估的不是需求，而是信用传导的时滞; [R015] 中国奢侈品市场的真正信号不是复苏，而是结构性分化加速; [R026] 物流成本正在重新定价中国消费股的利润表; [R021] 四月车市真正的信号不是销量下滑，而是库存去化正在重塑竞争格局; [R039] 中国房地产：反弹还是陷阱？市场低估了供给侧的再定价压力}

\section{区域市场与全球贸易}
\textbf{亚洲出口复苏由AI驱动的结构性分化主导，韩国半导体周期正从‘补库存’切换为‘再投资’，日本贸易数据揭示能源供给的‘价格-数量’错位。}

\begin{itemize}
  \item 亚洲出口领先指标NELI在6月攀升至117.7创2010年7月以来最高，但驱动力高度集中于半导体和AI相关链条，台湾和韩国显著跑赢区域平均水平，依赖大宗商品或低端制造的东南亚经济体受益有限。
  \item 韩国5月前20天日均出口环比增长2.8\%扭转颓势，半导体出口连续第8个月正增长且增速加快至11.9\%，半导体生产设备进口增速加速至21.3\%创去年9月以来最强，表明产业正从‘被动补货’转向‘主动扩产’。
  \item 日本四月贸易数据转正但出口量环比下降2.9\%，原油进口价格飙升57\%而进口量暴跌63.7\%，这一罕见组合指向全球供应链重构背景下日本能源安全逻辑正在被重新定价。
  \item 全球专业投资者的风险偏好框架正在经历结构性重置，47\%的受访者认为美股是2026年表现最佳的资产类别，看好欧洲股市的仅有7\%，AI资本支出被视为驱动市场表现的压倒性力量。
\end{itemize}
\begin{figure}[htbp]
\centering
\includegraphics[width=0.88\linewidth]{figures/fig_053.jpg}
\caption{Exhibit 1 - Exhibit 1: Rebound in early May exports}
\end{figure}
\begin{figure}[htbp]
\centering
\includegraphics[width=0.88\linewidth]{figures/fig_054.jpg}
\caption{Figure 1 - Figure 1: What will be the best performing asset class for the remainder of 2026?}
\end{figure}
\begin{figure}[htbp]
\centering
\includegraphics[width=0.88\linewidth]{figures/fig_048.jpg}
\caption{Exhibit 1 - Exhibit 1: Exports, Imports and Trade Balance}
\end{figure}
\begin{figure}[htbp]
\centering
\includegraphics[width=0.88\linewidth]{figures/fig_049.jpg}
\caption{Exhibit 2 - Exhibit 2: Export Volume by Region (Seasonally Adjusted Level)}
\end{figure}
\begin{figure}[htbp]
\centering
\includegraphics[width=0.88\linewidth]{figures/fig_050.jpg}
\caption{Exhibit 3 - Exhibit 3: Japan's Imported Crude Oil and LNG Volume}
\end{figure}
\begin{figure}[htbp]
\centering
\includegraphics[width=0.88\linewidth]{figures/fig_054.jpg}
\caption{Figure 1 - Figure 1: What will be the best performing asset class for the remainder of 2026?}
\end{figure}
\begin{figure}[htbp]
\centering
\includegraphics[width=0.88\linewidth]{figures/fig_055.jpg}
\caption{Figure 3 - Figure 3: How will Europe's defense spending surge reshape the investment landscape?}
\end{figure}
{\small\color{refgray}\textbf{References:} [R010] 亚洲出口复苏的真正信号：不再是普涨，而是AI驱动的结构性分化; [R007] 韩国出口反弹的真正信号：半导体周期正在从“补库存”切换为“再投资”; [R006] 日本贸易数据揭示的真正风险：不是出口放缓，而是能源供给的“价格-数量”错位; [R008] 市场真正低估的不是AI泡沫，而是资本从“恐惧地缘”到“恐惧踏空”的范式迁移}

\section{风险偏好与资金流向}
\textbf{主动管理基金的结构性错位与拥挤交易的二次筛选正在重塑市场定价逻辑，资本从‘恐惧地缘’转向‘恐惧踏空’的范式迁移是当前最被低估的力量。}

\begin{itemize}
  \item 全球专业投资者的‘存在性威胁’已从地缘政治切换为对AI踏空的恐惧，这一转变将重塑所有资产的定价锚，AI资本支出被视为‘自驱动’的盈利增长引擎。
  \item 主动基金正在以2012年以来从未有过的力度从软件股转向半导体股，但只有29\%的大型基金跑赢基准，远低于37\%的历史平均水平，主动管理能力困境正在加剧。
  \item 拥挤本身不是风险信号，拥挤叠加了错误的动量信号才是，过去12个月‘拥挤正向’组合跑赢‘拥挤负向’组合36个百分点，核心不是‘多少人买’而是‘买的人还在不在加码’。
  \item 全球主动型长线基金在4月净买入澳大利亚和印度，但净卖出中国和韩国，MSCI亚太（除日本）指数上涨15\%的背景下资金流向呈现高度分化。
  \item 美国电力需求到2050年需要翻倍，仅数据中心一项到2030年就将新增约120吉瓦需求，没有任何单一技术路径能够独自满足，‘所有技术都供不应求’的格局正在重新定义基础设施投资周期。
\end{itemize}
\begin{figure}[htbp]
\centering
\includegraphics[width=0.88\linewidth]{figures/fig_054.jpg}
\caption{Figure 1 - Figure 1: What will be the best performing asset class for the remainder of 2026?}
\end{figure}
\begin{figure}[htbp]
\centering
\includegraphics[width=0.88\linewidth]{figures/fig_283.jpg}
\caption{Exhibit 6 - Exhibit 6: Mutual funds added in Semis, cut in Software excludes AAPL, AVGO, MSFT, NVDA}
\end{figure}
\begin{figure}[htbp]
\centering
\includegraphics[width=0.88\linewidth]{figures/fig_284.jpg}
\caption{Exhibit 7 - Exhibit 7: Mutual fund tilt in Semis vs. Software ex-mega caps is largest since 2012}
\end{figure}
\begin{figure}[htbp]
\centering
\includegraphics[width=0.88\linewidth]{figures/fig_286.jpg}
\caption{Exhibit 11 - Exhibit 11: Mutual funds were 723 bp underweight the Mag 7}
\end{figure}
\begin{figure}[htbp]
\centering
\includegraphics[width=0.88\linewidth]{figures/fig_279.jpg}
\caption{Exhibit 23 - Exhibit 1: 29\% of large-cap mutual funds are outperforming their benchmarks}
\end{figure}
\begin{figure}[htbp]
\centering
\includegraphics[width=0.88\linewidth]{figures/fig_280.jpg}
\caption{Exhibit 3 - Exhibit 3: Mutual fund cash balances have risen YTD as of March 31, 2026}
\end{figure}
{\small\color{refgray}\textbf{References:} [R008] 市场真正低估的不是AI泡沫，而是资本从“恐惧地缘”到“恐惧踏空”的范式迁移; [R025] 市场真正低估的不是AI需求，而是主动管理基金的结构性错位; [R023] 拥挤交易并非风险信号，真正的超额收益来自对拥挤度的二次筛选; [R033] 美国电力需求的真实拐点：不是“更多新能源”，而是“所有技术都供不应求”; [R017] 数据中心投资的下一个分水岭不是算力，而是“表后电力”的规模化落地}

\section*{结语}
当前市场正处于多重结构性力量交汇的节点：AI与机器人主导的技术周期正在重塑增长方程式的分母，而宏观层面的加息风险与权益风险溢价压缩则为资产定价设定了新的约束边界。理解这些结构性分化背后的驱动力，远比追逐短期数据波动更为重要。
\vfill
{\small\color{refgray}Personal reading notes and learning share only. Not investment advice.}
\end{document}
